\section{The alternating Gram distribution}\label{sec:gram}

The additive group $\Alt_a(\F_q)$ of alternating $a$ by $a$ matrices
has size $q^D$, where $D=\binom a2$.
For $v\ge0$, take $Y_v$ uniform on $\Mat_{2v,a}(\F_q)$ and define
\begin{equation}\label{eq:Gamma}
 \Gamma_v=Y_v^{\tr}J_vY_v,\qquad
 J_v=\begin{pmatrix}0&I_v\\-I_v&0\end{pmatrix}.
\end{equation}
When $v=0$ this is the zero matrix.

\begin{lemma}[Exact Fourier transform]\label{lem:fourier}
Fix a nontrivial additive character $\psi$ of $\F_q$ and use
upper-triangular coordinates to identify the dual of $\Alt_a$.
The character indexed by $Z\in\Alt_a$ has expectation
\begin{equation}\label{eq:fourier}
 \widehat{\mathcal L(\Gamma_v)}(Z)=q^{-v\rank Z}.
\end{equation}
For $a\ge2$, this implies, for every $T\in\Alt_a$,
\begin{equation}\label{eq:gramuniform}
 \abs{q^D\PP(\Gamma_v=T)-1}\le\varepsilon_v,\qquad
 \varepsilon_v=(q^D-1)q^{-2v}.
\end{equation}
\end{lemma}

\begin{proof}
Group the rows of $Y_v$ into independent uniform pairs $(x_i,y_i)$.
The Gram matrix is the sum of their alternating products. Pairing
one product with the upper-triangular coordinates of $Z$ gives
$x_i^{\tr}Zy_i$, also in characteristic two.
For fixed $x_i$, averaging the character over $y_i$ gives one if
$x_i^{\tr}Z=0$ and zero otherwise. Averaging over $x_i$ gives
$q^{-\rank Z}$. Independence proves \eqref{eq:fourier}.
Every nonzero alternating matrix has rank at least two, so
Fourier inversion yields \eqref{eq:gramuniform}.
\end{proof}

For $0\le r\le\lfloor a/2\rfloor$, define
\begin{align}
 \mathcal Z_c(a,r)
   &=\sum_{k=0}^{a-2r}\gb{a-2r}{k}\,
                    \mathcal E_c(a-k,r),\label{eq:Z}\\
 \mathcal A_a(r)
   &=\gb a{2r}\frac{|\GL_{2r}(\F_q)|}{|\Sp_{2r}(\F_q)|}.
                    \label{eq:A}
\end{align}
Both group orders in the ratio are one when $r=0$.
The first number counts all linear maps to a symplectic $2c$-space
with a specified Gram matrix of rank $2r$, allowing a kernel:
choose any $k$-subspace of the radical as kernel, then embed the quotient.
The second counts alternating matrices of rank $2r$: choose a radical,
then a nondegenerate alternating form on the quotient, with
$|\GL_{2r}|/|\Sp_{2r}|$ choices.

Define the \emph{complement discrepancy}
\begin{equation}\label{eq:Delta}
 \Delta_{q,a,c}
 =\frac12\sum_{r=0}^{\lfloor a/2\rfloor}
 \mathcal A_a(r)
 \left|\frac{\mathcal Z_c(a,r)}{q^{2ca}}-q^{-D}\right|.
\end{equation}
It equals $\TV(\Gamma_c,U_{\Alt})$, where $U_{\Alt}$ is uniform on
$\Alt_a(\F_q)$. Thus \eqref{eq:Delta} is an explicit finite formula
with at most $\lfloor a/2\rfloor+1$ absolute values.

\section{A universal boundary law and rank blindness}\label{sec:boundary}

Fix $q$, $a\ge2$, an alternating $H$ of size $a$ and rank $2h$, and
$c\ge0$. For each sufficiently large $s$, with $s+c\ge a-h$, decompose
\[
 V_s=U_s\perp W,\qquad \dim U_s=2s,\quad\dim W=2c.
\]
The displayed inequality ensures that an injective source frame
with Gram matrix $H$ exists.
Let $X_s:\F_q^a\to V_s$ be uniform among such frames, and let $A_s$
be its projection onto $U_s$ in a symplectic basis.
Let $Z_s$ be a uniform $2s$ by $a$ matrix.
Multiplication by the target Gram matrix converts this projection
into \eqref{eq:compression}, without changing rank or distance from
uniform matrices.

\begin{theorem}[Complement-Gram boundary law]\label{thm:boundary}
With $q,a,c,H$ fixed,
\begin{align}
 \TV(A_s,Z_s)&=\Delta_{q,a,c}+O_{q,a,c}(q^{-2s}),\label{eq:boundary}\\
 \TV(\rank A_s,\rank Z_s)&=O_{q,a,c}(q^{-2s}).\label{eq:rankblind}
\end{align}
The full-matrix limit is independent of the source form $H$, including
its rank. Moreover,
\begin{equation}\label{eq:Deltabounds}
 \frac12q^{-2c}\le\Delta_{q,a,c}
 \le\min\left\{1,\frac12\sqrt{q^D-1}\,q^{-2c}\right\}.
\end{equation}
It is strictly positive for fixed finite $c$, and for fixed $q,a$
its decay as $c$ grows is $\Theta_{q,a}(q^{-2c})$.
\end{theorem}

A fixed-dimensional hidden symplectic complement therefore leaves
a persistent constraint on the full observed matrix, although all
rank distributions approach the unrestricted one. The surviving
constraint is in the alternating Gram matrix.

\begin{proof}
Put $m=s+c$ and start with unrestricted independent uniform blocks
$X=(A;Y)$, with $2s$ observed and $2c$ hidden symplectic coordinates.
The Gram matrix is $\Gamma_s+\Gamma_c$.
Let $E$ denote the event that it equals $H$, and write
\[
 F_c(T)=q^D\PP(\Gamma_c=T),\qquad
 R_m=q^D\PP(\Gamma_m=H).
\]
The conditional observed law $Q_s=\mathcal L(A\mid E)$ has
likelihood, relative to uniform observed matrices,
\begin{equation}\label{eq:conditionaldensity}
 \frac{\dd Q_s}{\dd\mathcal L(Z_s)}(A)
 =\frac{F_c(H-A^{\tr}J_sA)}{R_m}.
\end{equation}
In particular, $\E F_c(H-\Gamma_s)=R_m$.

For sufficiently large $s$, $\varepsilon_m<1$.
By Lemma \ref{lem:fourier}, $|R_m-1|\le\varepsilon_m$.
Set $f(T)=|F_c(H-T)-1|$. Translation and negation preserve uniform
measure on $\Alt_a$, so $\E_{U_{\Alt}}f=2\Delta_{q,a,c}$.
The pointwise relative bound \eqref{eq:gramuniform}, together with
$f\ge0$, yields
\[
 \abs{\E_{\Gamma_s}f-2\Delta_{q,a,c}}
 \le2\varepsilon_s\Delta_{q,a,c}.
\]
Replacing $F_c/R_m$ by $F_c$ inside the absolute value in total
variation costs at most
$\tfrac12|R_m^{-1}-1|\E_{\Gamma_s}F_c
=\tfrac12|R_m-1|$.
Hence
\begin{equation}\label{eq:Qbound}
 \abs{\TV(Q_s,Z_s)-\Delta_{q,a,c}}
 \le\tfrac12\varepsilon_m+\varepsilon_s\Delta_{q,a,c}.
\end{equation}

To pass to injective source frames, let
$L_a=(q^a-1)/(q-1)$ be the number of lines in $\F_q^a$.
An unrestricted $2m$ by $a$ matrix kills a given line with probability
$q^{-2m}$. The union bound gives
\[
 \PP(X\text{ not injective}\mid E)
 \le\eta_m:=\frac{q^DL_aq^{-2m}}{1-\varepsilon_m}.
\]
Conditioning on injectivity changes the joint law by the probability
discarded, and projection cannot increase total variation. The
conditioned law is precisely the uniform injective frame with Gram $H$.
Therefore
\begin{equation}\label{eq:expliciterror}
 \abs{\TV(A_s,Z_s)-\Delta_{q,a,c}}
 \le\eta_m+\tfrac12\varepsilon_m+
                   \varepsilon_s\Delta_{q,a,c}.
\end{equation}
This proves \eqref{eq:boundary} uniformly over all source forms of
the fixed size. When $H$ is nonsingular, its Gram condition already
forces injectivity, so the term $\eta_m$ is unnecessary.

For the rank estimate, rewrite the source-frame count as
\begin{equation}\label{eq:normalizedE}
 \mathcal E_m(a,h)=q^{2ma-D}C_{m,a,h},\qquad
 C_{m,a,h}=\prod_{j=m-a+h+1}^{m}(1-q^{-2j}).
\end{equation}
On the feasible range this product is at least
$\prod_{j\ge1}(1-q^{-2j})>0$.
The uniform observed matrix has probability at most $L_aq^{-2s}$
of failing full column rank. Dividing by the probability
$q^{-D}C_{m,a,h}$ of the injective Gram event gives
\[
 \PP(\rank A_s<a)\le
 \frac{q^D}{C_{m,a,h}}L_aq^{-2s}.
\]
Comparing both rank distributions with the point mass at $a$ yields
the explicit bound
\begin{equation}\label{eq:rankexplicit}
 \TV(\rank A_s,\rank Z_s)
 \le\left(1+\frac{q^D}{C_{m,a,h}}\right)L_aq^{-2s}.
\end{equation}

Finally, a character indexed by a rank-two alternating matrix has
expectation $q^{-2c}$ under $\Gamma_c$ and zero under uniform measure.
Expectations of a complex function of modulus at most one differ by
at most twice total variation, proving the lower bound in
\eqref{eq:Deltabounds}. Parseval and Cauchy--Schwarz give
\[
 \Delta_{q,a,c}
 \le\frac12\left(\sum_{T\ne0}q^{-2c\rank T}\right)^{1/2}
 \le\frac12\sqrt{q^D-1}\,q^{-2c}.
\]
The trivial upper bound one completes the proof.
\end{proof}

\begin{corollary}[Small source dimensions]\label{cor:small}
For $a=2$ or $3$,
\[
 \Delta_{q,a,c}=(1-q^{-\binom a2})q^{-2c}.
\]
\end{corollary}

\begin{proof}
Every nonzero dual alternating matrix in these dimensions has rank two.
Thus \eqref{eq:fourier} identifies the Gram law as
$q^{-2c}\delta_0+(1-q^{-2c})U_{\Alt}$. Its distance from uniform is
the displayed quantity.
\end{proof}

\subsection{An exact likelihood for a nonsingular source}

If $a=2d$ and $H$ is nonsingular, \eqref{eq:conditionaldensity}
already gives the actual projected-frame law:
\begin{equation}\label{eq:nondegdensity}
 \frac{\PP(A_s=A)}{q^{-2sa}}=
 \frac{q^D\PP(\Gamma_c=H-A^{\tr}J_sA)}
      {\prod_{j=s+c-d+1}^{s+c}(1-q^{-2j})}.
\end{equation}
The hidden complement must supply exactly the missing Gram matrix.
The support criterion becomes
$\rank(H-A^{\tr}J_sA)\le2c$.
It is sufficient: embed the nondegenerate part of this missing form
and send its radical to zero. Nonsingularity of $H$ guarantees that
the combined observed and hidden map is injective.

\section{The plane case: exact optimal tests}\label{sec:plane}

Take $a=2$ with source form
$H=\left(\begin{smallmatrix}0&1\\-1&0\end{smallmatrix}\right)$.
Let $s\ge1$, $c\ge0$ and $A=(u_1,u_2)$ be the observed projection
onto a nondegenerate $2s$-space. Set
\[
 x=q^{-2s},\quad y=q^{-2c},\quad
 \tau(A)=u_1^{\tr}J_su_2,
\]
and let $Z$ be a uniform matrix of the same size.

\begin{theorem}[Plane law and rank information loss]\label{thm:plane}
The exact likelihood is
\begin{equation}\label{eq:planeL}
 \frac{\PP(A=C)}{q^{-4s}}=
 \begin{cases}
 \displaystyle\frac{1+(q-1)y}{1-xy},&\tau(C)=1,\\[6pt]
 \displaystyle\frac{1-y}{1-xy},&\tau(C)\ne1.
 \end{cases}
\end{equation}
The event $\tau(C)=1$ is an optimal distinguishing event, and
\begin{equation}\label{eq:planeTV}
 \boxed{\TV(A,Z)=\frac{y(1-x)(q-1+x)}{q(1-xy)}.}
\end{equation}
The optimal rank-only event is full column rank, with
\begin{equation}\label{eq:planeRank}
 \boxed{\TV(\rank A,\rank Z)=
 \frac{yx(1-x)(q+1-qx)}{1-xy}.}
\end{equation}
For fixed $q,c$, as $s\to\infty$,
\[
 \TV(A,Z)\longrightarrow(1-q^{-1})q^{-2c},
 \qquad
 \TV(\rank A,\rank Z)\sim(q+1)q^{-2(s+c)}.
\]
Both source and target are nondegenerate; neither is Lagrangian.
\end{theorem}

\begin{proof}
In a symplectic $2c$-space, the number of ordered pairs with given
nonzero pairing is $(q^{2c}-1)q^{2c-1}$.
For pairing zero the number is
$q^{2c}+(q^{2c}-1)q^{2c-1}$: either the first vector is zero,
or its pairing imposes one linear equation on the second.
At $c=0$ these expressions are interpreted as zero and one.
The hidden pairing must be $1-\tau(C)$. Divide by the total number
$(q^{2(s+c)}-1)q^{2(s+c)-1}$ of symplectic pairs in the ambient
space to obtain \eqref{eq:planeL}.

For a uniform observed matrix, $\PP(\tau(Z)=1)=q^{-1}(1-x)$.
The first likelihood exceeds one, and the second is below one.
Thus total variation is the excess mass of the pairing-one event,
which gives \eqref{eq:planeTV}.

A rank-deficient two-column matrix has pairing zero and hence the
same lower likelihood. Its uniform probability is
\[
 D_s=1-(1-x)(1-qx)=(q+1)x-qx^2.
\]
All deficits in the rank law are at ranks zero and one, and the
compensating excess is at rank two. The rank distance is
$D_s[1-(1-y)/(1-xy)]$, giving \eqref{eq:planeRank}.
The asymptotics follow by letting $x\to0$.
\end{proof}

\begin{example}[Twenty observed binary coordinates]\label{ex:binary}
For $q=2,s=10,c=1$, the ambient space has dimension $22$, the target
has dimension $20$, and the source has dimension $2$. The formulas give
\[
 \TV(A,Z)=0.1250000298022158\ldots,\qquad
 \TV(\rank A,\rank Z)=7.152547709667702\ldots\times10^{-7}.
\]
The pairing test retains a probability discrepancy of about one eighth,
while the best rank-only test has discrepancy below one millionth.
These are differences between event probabilities, not biases of
individual matrix entries.
\end{example}
